\begin{helper}[Strict augmentation beyond the global order]
Fix the Scheide configuration and replacement family.  Suppose there
exist $\alpha,d,f,w,u$ with the following package.  Put
\[
 \varphi_d=\noderef{TwoColorComponentSwap}
 (G,\varphi,\gamma(d),d,\operatorname{inl}v^0_d),
\]
\[
 Q=\noderef{TwoColorComponent}
 (G,\varphi_d,\alpha_0,\alpha,\operatorname{inl}u),\qquad
 \psi=\noderef{TwoColorComponentSwap}
 (G,\varphi_d,\alpha_0,\alpha,\operatorname{inl}u),
\]
and $T_d=(R d).\mathit{replacementTree}$.  The package is:
$\alpha\ne\alpha_0$; $G.\mathop{\rm ends}(f)=s(w,u)$;
$\varphi_d(f)=\alpha$; $\operatorname{inr}f\in Q$; no vertex-node of
$V(T)$ lies in $Q$; $(T_d,e_0,\psi)$ is globally maximum;
$\alpha_0$ is missing at $y_0=C.y(0)$ under $\varphi_d$;
$y_0,w\in V(T_d)$; $u\notin V(T_d)$; and $f$ is not an edge identity of
$T_d$.  Then there are a partial $k$-edge-coloring $\psi'$ and a
tree-sequence $S$ such that $(G,e_0,\psi')$ is a $k$-triple, $S$ is a
Tashkinov tree for this triple, and
\[
 |V(S)|=\noderef{TashkinovOrder}(G)+1.
\]
\end{helper}
\begin{proof}
Choose the displayed package.  Since
$\operatorname{inr}f\in Q$, $\varphi_d(f)=\alpha$, and
$\alpha\ne\alpha_0$, unfolding
\noderef{TwoColorComponentSwap} gives
\[
 \psi(f)=\alpha_0.                                        \tag{1}
\]

The base-position equation in \noderef{ScheideStructuralConfiguration}
puts $y_0$ in $V(T)$.  Hence the avoidance assumption gives
$\operatorname{inl}y_0\notin Q$.  No edge-node of $Q$ is incident with
$y_0$: otherwise the adjacency clause in
\noderef{TwoColorIncidenceGraph} would put $\operatorname{inl}y_0$ in
the same connected component.  The second switch therefore fixes every
edge incident with $y_0$, so it fixes the whole set
\noderef{MissingColors}$(G,\varphi_d,y_0)$.  The packaged missingness
thus gives
\[
 \alpha_0\in\noderef{MissingColors}(G,\psi,y_0).           \tag{2}
\]

The endpoint equation, $w\in V(T_d)$, and $u\notin V(T_d)$ put $f$ in
\noderef{BoundaryEdges}$(G,V(T_d))$.  Apply
\noderef{TreeSequenceBoundaryEdgeAppendExists} to obtain a tree-sequence
$S$ with \noderef{TreeSequenceAdjoinEdge}$(T_d,f,S)$.  Unfolding the
latter relation says that $S$ has length $T_d.length+1$, keeps every old
vertex and edge position, and has $f$ as its final edge.  Its final new
vertex is the outside end $u$ by the attachment equation and vertex
injectivity.  More explicitly, $u\notin V(T_d)$ and the hypothesis that
$f$ is not an old edge identity are exactly the freshness facts needed by
the vertex- and edge-injectivity clauses of \noderef{TreeSequence}; the
boundary-append lemma packages these facts into the asserted tree-sequence.

The globally maximum package gives
\noderef{TashkinovTree}$(G,e_0,\psi,T_d)$.  Every old noninitial edge of
$S$ has exactly the same earlier missing witness as in $T_d$, because
the initial-segment equations preserve both positions.  For the new last
edge, (1) gives its color $\alpha_0$, while (2) gives a missing witness at
the earlier vertex $y_0\in V(T_d)$.  Thus every noninitial position of
$S$ satisfies the Tashkinov condition, and $S$ is a Tashkinov tree for
the same $k$-triple.

By \noderef{TreeSequenceVertexSet} and the vertex-injectivity clause of
\noderef{TreeSequence}, the displayed vertex map of a tree-sequence of
length $m$ is an injection from exactly $m+1$ indices onto its vertex set.
The adjoin equation therefore gives
\[
 |V(S)|=|V(T_d)|+1.                                       \tag{3}
\]
The cardinality clause of \noderef{GlobalMaximumTriple} says
$|V(T_d)|=\noderef{TashkinovOrder}(G)$.  Its
\noderef{KTriple} clause supplies the same $k$-triple for the unchanged
coloring $\psi$.  Taking $\psi'=\psi$, equation (3) is therefore
$|V(S)|=\noderef{TashkinovOrder}(G)+1$, and the constructed $S$ has all
the properties asserted in the conclusion.
\end{proof}
